%%%%%%%%%%%%%%%%%%%%%%%%%%%%%%%%%%%%%%%%%%%%%%%%%%%%%%%
%%                                                   %%
%% Modelo LaTeX NÃO OFICIAL para trabalhos           %%
%% acadêmicos da UFERSA, mantido por                 %%
%% Enthony Araujo.                                   %%
%%                                                   %%
%% Versão modificada do modelo da UFC (Ednardo       %%
%% Moreira Rodrigues, Alan Batista de Oliveira e     %%
%% colaboradores), adaptado do ueceTeX2 (Thiago      %%
%% Nascimento, https://github.com/thiagodnf/uecetex2)%%
%% e baseado no abnTeX2. Créditos completos e        %%
%% licença (LPPL 1.3c) em lib/ufersatex.sty,         %%
%% README.md e LICENSE.                              %%
%%                                                   %%
%% Instruções de uso: README.md                      %%
%%                                                   %%
%%%%%%%%%%%%%%%%%%%%%%%%%%%%%%%%%%%%%%%%%%%%%%%%%%%%%%%

% Importações de pacotes
\usepackage[utf8]{inputenc}                         % Acentuação direta
\usepackage[T1]{fontenc}                            % Codificação da fonte em 8 bits
\usepackage{graphicx}                               % Inserir figuras
\usepackage{amsmath,amsthm,amsfonts,amssymb}        % Fontes e símbolos matemáticos
\usepackage{booktabs}                               % Comandos para tabelas
\usepackage{verbatim}                               % Texto é interpretado como escrito no documento
\usepackage{multirow, array}                        % Múltiplas linhas e colunas em tabelas
\usepackage{indentfirst}                            % Endenta o primeiro parágrafo de cada seção.
\usepackage{listings}                               % Utilizar codigo fonte no documento
\usepackage{xcolor}                                 % Cores
\usepackage{microtype}                              % Melhorias de justificação
\usepackage[portuguese,ruled,lined]{algorithm2e}    % Escrever algoritmos
\usepackage{tocloft}                                % Permite alterar a formatação do Sumário
\usepackage{etoolbox}                               % Usado para alterar a fonte da Section no Sumário
\usepackage{enumitem}                               % Listas (alíneas)
\usepackage[nogroupskip,nonumberlist]{glossaries}   % Lista de abreviaturas e siglas e glossário
\usepackage{caption}                                % Altera o comportamento da tag caption

% Citações e referências no sistema autor-data, conforme NBR 10520:2023 e
% NBR 6023:2018:
%  - \citeoption{ufersa-nbr10520-2023} (abaixo): sobrenome em maiúsculas e
%    minúsculas também dentro dos parênteses, ex.: (Silva, 2023);
%    ver lib/abnt-opcoes-ufersa.bib;
%  - abnt-etal-cite=3: "et al." na citação para quatro ou mais autores;
%  - abnt-etal-list=0: todos os autores na lista de referências;
%  - abnt-emphasize=bf: destaque do título em negrito.
\usepackage[alf, abnt-emphasize=bf, recuo=0cm, abnt-etal-cite=3, abnt-etal-list=0, abnt-etal-text=it]{abntex2cite}
\AtBeginDocument{\citeoption{ufersa-nbr10520-2023}}

\usepackage{chngcntr}
\counterwithout{equation}{chapter}                  % Equações numeradas (1), (2), ...

%%%%%%%%%%%%%%%%%%% FONTE DO TEXTO %%%%%%%%%%%%%%%%%%%%%%%%%%%%%%%
% O template da UFERSA aceita Times ou Arial, tamanho 12.
% Padrão: Times New Roman.
\usepackage{mathptmx}         % Usa a fonte Times New Roman
% Para usar Arial, comente a linha acima e descomente as duas abaixo:
%\usepackage{helvet}
%\renewcommand{\familydefault}{\sfdefault}
%%%%%%%%%%%%%%%%%%%%%%%%%%%%%%%%%%%%%%%%%%%%%%%%%%%%%%%%%%%%%%%%%%

\usepackage{pdfpages}         % Incluir pdf no documento (ficha catalográfica, errata etc.)
\usepackage{lib/ufersatex}    % Configurações do modelo (UFERSA)

\makeglossaries % Organiza e gera a lista de abreviaturas, simbolos e glossario
\makeindex      % Gera o Indice do documento
